\section{Nuclear averaging and bounded conditional moments}
\label{u:sec:averaging}

Throughout this section, $A$ is simple, separable, unital,
infinite dimensional, and nuclear, has strict comparison as in
Definition~\ref{u:def:strict-comparison}, and satisfies
$T(A)\ne\varnothing$.  Retain $D$ and $\Lambda$ from
Section~\ref{u:sec:preliminaries}.  Nuclearity will provide
finite-dimensional approximations whose errors are small on average.
The centered symmetries of Proposition~\ref{u:prop:centered-symmetry}
then turn those averages into one selfadjoint observable.  A fourth-moment
estimate permits a norm cutoff with controlled loss of its first two
coefficient moments.  The symmetries used here need not commute with
one another.  The bounded local moment estimates obtained after the
cutoff will allow us to construct commuting variables in the next
section, by enlarging the coefficient algebra at each step.

\subsection{Convex approximations and homomorphism replacement}

A completely positive map has \emph{order zero} if it preserves
orthogonality of positive elements: $xy=0$ for $x,y\ge0$ implies
$\phi(x)\phi(y)=0$.  We use the following external approximation theorem
in its full nuclear scope.

\begin{theorem}[Hirshberg--Kirchberg--White, {\cite[Theorem~1.4]{HKW}}]
\label{u:thm:hkw-approximation}
Let $N$ be a nuclear $C^*$-algebra, let $K\subset N$ be finite, and let
$\varepsilon>0$.  There are a finite-dimensional $C^*$-algebra $F$ and
completely positive contractive maps
\[
 N\xrightarrow{\psi}F\xrightarrow{\phi}N
\]
such that $\norm{\phi\psi(x)-x}<\varepsilon$ for $x\in K$, and
\[
 \phi=\sum_{i=1}^{M}\beta_i\phi_i,
 \qquad \beta_i\ge0,\qquad \sum_{i=1}^{M}\beta_i=1,
\]
where each $\phi_i:F\to N$ is completely positive, contractive, and of
order zero.  The number $M$ is finite for the chosen approximation;
no bound on $M$ independent of $K$ and $\varepsilon$ is asserted.
\end{theorem}

In particular, we may write an approximating map in the form
$\sum_i\beta_i\phi_i\psi_i$ with finite-dimensional $F_i$ and completely
positive contractive maps $N\xrightarrow{\psi_i}F_i
\xrightarrow{\phi_i}N$, each $\phi_i$ of order zero: simply take
$F_i=F$ and $\psi_i=\psi$.  This theorem does not assume a bound on nuclear
dimension.  Its unrestricted number of summands is the reason for the
cutoff at the end of this section.

\begin{lemma}[Homomorphism replacement]
\label{u:lem:homomorphism-replacement}
Let $F$ be finite dimensional, let $\phi:F\to D$ be a completely positive
contractive order-zero map, and put $h=\phi(1_F)$.  For $0<\eta<1$ there
is a possibly nonunital $*$-homomorphism $\rho:F\to D$ such that, writing
$P=\rho(1_F)$, for every $\tau\in\Lambda$,
\begin{equation}
 \tau(P)\ge\tau(h)-\eta,
 \qquad
 \norm{\phi(c)-\rho(c)}_{2,\tau}^{2}
 \le 4(1-\tau(h))+2\eta
 \quad (c\in F,\ \norm{c}\le1).
 \label{u:eq:homomorphism-replacement}
\end{equation}
\end{lemma}

\begin{proof}
We first construct a supporting homomorphism $\pi$ in $D^{**}$, then
use the supplied projections to transport its matrix units into $D$.
The final estimates measure the total trace lost in this transport,
independently of the matrix sizes.  The supporting-homomorphism step
is the finite-dimensional case of the order-zero structure theorem
\cite[Theorem~3.3]{WZ}; we include the details to specify the support
used below.  Put $q=\supp(h)\in D^{**}$.  For every projection $e\in F$,
the positive elements $\phi(e)$ and $\phi(1_F-e)$ are orthogonal.
Consequently $\phi(e)$ commutes with their sum $h$.  Projections linearly
span $F$, so $h$ commutes with $\phi(F)$.  For $t>0$ define
\[
 \pi_t(x)=(h+t1)^{-1/2}\phi(x)(h+t1)^{-1/2}
         =(h+t1)^{-1}\phi(x).
\]
This map is completely positive and contractive, since
$0\le\pi_t(1_F)=h(h+t1)^{-1}\le1$.  If $e$ is a projection, orthogonality
of $\phi(e)$ and $\phi(1_F-e)$ gives
\[
 \pi_t(e)=\phi(e)(\phi(e)+t1)^{-1}
 \ \longrightarrow\ \supp(\phi(e))
 \quad\hbox{strongly as }t\downarrow0.
\]
Because projections span $F$, $\pi_t$ has a pointwise strong limit
$\pi:F\to D^{**}$.  Positivity at every matrix level passes to this
limit.  Thus $\pi$ is completely positive, $\pi(1_F)=q$, and $\pi$ takes
every projection to a projection.  Its range lies in $qD^{**}q$:
indeed, for positive $x\in F$ one has
$0\le\pi(x)\le\norm{x}q$, and then one uses linearity.
Viewed as a unital completely positive map into that corner, $\pi$ is
multiplicative.  To recall the multiplicative-domain argument, equality
in the Schwarz inequality at a projection $e$ forces
$\pi(ex)=\pi(e)\pi(x)$ and $\pi(xe)=\pi(x)\pi(e)$ for every $x$;
this follows by applying Schwarz to $e+zx$ for all $z\in\C$ and comparing
the linear terms.  Applying this to all projections and using their
linear span proves multiplicativity.  Finally,
\[
 h\pi_t(x)=\phi(x)-t\pi_t(x)\longrightarrow\phi(x)
 \quad\hbox{in norm},
\]
so $h\pi(x)=\phi(x)$.  The commutation with $h$ passes to the strong
limit, and hence
\begin{equation}
 \phi(x)=h\pi(x)=\pi(x)h,
 \qquad \pi(1_F)=\supp(h)=q.
 \label{u:eq:order-zero-support}
\end{equation}
The zero map is included, with $q=0$.

Write $F=\bigoplus_{\alpha}M_{n_\alpha}$, and choose matrix units
$e^{(\alpha)}_{st}$.  Work first in one block and omit its superscript.
Set $a=\phi(e_{11})$.  Proposition~\ref{u:prop:projection-supply} supplies
a projection $p\in D$ with
\[
 p\le\mathbf1_{[\eta,1]}(a)\le\pi(e_{11}),
 \qquad \tau(p)=\tau((a-\eta)_+)
 \quad(\tau\in\Lambda).
\]
The final support inequality follows from
$\pi(e_{11})=\supp(\phi(e_{11}))$.  On $[0,1]$ define
\[
 k(s)=\min(s/\eta,1),\qquad
 g(s)=
 \begin{cases}k(s)/s,&s>0,\\1/\eta,&s=0.\end{cases}
\]
Both functions are continuous, and $k(s)=s g(s)$.  The spectral
containment of $p$ gives $k(a)p=pk(a)=p$.  In particular the elements
\begin{equation}
 v_s=\pi(e_{s1})k(a)=\phi(e_{s1})g(a)\in D
 \label{u:eq:replacement-membership}
\end{equation}
are in $D$, even though the individual $\pi(e_{s1})$ were constructed
only in $D^{**}$.  The equality uses
$\pi(e_{s1})a=\phi(e_{s1})$, which follows from
Equation~\eqref{u:eq:order-zero-support}.  Define
\[
 \rho(e_{st})=v_s p v_t^*
             =\pi(e_{s1})p\pi(e_{1t})\in D.
\]
Since $p\le\pi(e_{11})$, multiplication gives
\[
 \rho(e_{st})\rho(e_{uv})=\delta_{tu}\rho(e_{sv}),
 \qquad \rho(e_{st})^*=\rho(e_{ts}).
\]
For different blocks the products vanish, because their supporting
matrix units are orthogonal.  Thus these formulas extend to a
$*$-homomorphism $\rho:F\to D$.  Direct multiplication by the matrix
units also shows that, for $P=\rho(1_F)$,
\[
 P\le q,\qquad P\pi(x)=\pi(x)P=\rho(x)\quad(x\in F).
\]
This construction uses the support containment of $p$; it does not
require $p$ to commute with $a$, nor $P$ to commute with $h$.

Fix $\tau\in\Lambda$.  In all expressions on $D^{**}$ use its canonical
normal tracial extension, still denoted by $\tau$.  Faithfulness of this
extension is not needed.  Scalar functional calculus gives, in each
block,
\[
 0\le a-(a-\eta)_+\le\eta\pi(e_{11}),
 \qquad
 0\le\tau(a)-\tau(p)\le\eta\tau(\pi(e_{11})).
\]
Traciality and the partial isometries $\pi(e_{s1})$ imply that
$\phi(e_{ss})$, $\rho(e_{ss})$, and $\pi(e_{ss})$ have the same traces as
$a$, $p$, and $\pi(e_{11})$, respectively.  Summing over all diagonal
entries of all blocks therefore gives
\begin{equation}
 0\le\tau(h)-\tau(P)\le\eta\tau(q)\le\eta.
 \label{u:eq:replacement-unit-loss}
\end{equation}
The total loss is measured by the trace of $q$, so no matrix-size factor
appears.

For a contraction $c\in F$, Equation~\eqref{u:eq:order-zero-support} and
commutation with $\pi(c)$ give
\begin{align*}
 \norm{\phi(c)-\pi(c)}_{2,\tau}^2
 &\le\tau((q-h)^2)\le\tau(q-h)\le1-\tau(h),\\
 \norm{\rho(c)-\pi(c)}_{2,\tau}^2
 &\le\tau(q-P)\le1-\tau(h)+\eta.
\end{align*}
Here $q-h$ is a positive contraction and $q-P$ is a projection; both
commute with $\pi(c)$.  The squared triangle inequality now yields
Equation~\eqref{u:eq:homomorphism-replacement}.
\end{proof}

\subsection{Averaging the centered symmetries}

\begin{proposition}[An observable with exact coefficient moments]
\label{u:prop:averaged-observable}
Let $E\subset A$ be a finite set of selfadjoint contractions and let
$0<\eta<1$.  There is a selfadjoint $z\in D$ such that, for every
$\tau\in\Lambda$,
\begin{equation}
 \begin{aligned}
  \norm{[z,a]}_{2,\tau}^{2}&\le72\eta &&(a\in E),\\
  \tau(bz)&=0,\qquad \tau(bz^2)=\tau(b) &&(b\in A),\\
  \tau(z^4)&\le3.
 \end{aligned}
 \label{u:eq:averaged-observable}
\end{equation}
The construction gives
$z=\sum_{i=1}^M\sqrt{\beta_i}\,r_i$, where the $r_i$ are selfadjoint
unitaries, $\beta_i\ge0$, and $\sum_i\beta_i=1$.  In fact,
$\tau(z^4)\le3-2\sum_i\beta_i^2$.  There is no asserted norm bound on
$z$ independent of $M$.
\end{proposition}

\begin{proof}
Apply Theorem~\ref{u:thm:hkw-approximation} to the finite set
\[
 \{1_A\}\cup E\cup\{a^2:a\in E\}
\]
with error $\eta$.  In the finite-family notation following that theorem,
put $T_i=\phi_i\psi_i$ and $T_0=\sum_i\beta_iT_i$.  Thus
$\norm{T_0(x)-x}<\eta$ on this set.  Regard the outgoing maps as taking
values in $D$ and write $h_i=\phi_i(1_{F_i})$.  Fix one trace
$\tau\in\Lambda$ throughout the following weighted estimates.  Positivity
and contractivity give $T_i(1_A)\le h_i\le1$, and hence
\begin{equation}
 \sum_i\beta_i(1-\tau(h_i))
 \le\tau(1-T_0(1_A))<\eta.
 \label{u:eq:averaged-unit-defect}
\end{equation}
For $a\in E$, the Schwarz inequality for the completely positive
contraction $T_i$ gives $T_i(a)^2\le T_i(a^2)$.  Expansion of the squares
and traciality imply
\begin{align}
 \sum_i\beta_i\norm{a-T_i(a)}_{2,\tau}^{2}
 &\le\tau(a^2)-2\tau(aT_0(a))+\tau(T_0(a^2))\notag\\
 &= -2\tau(a(T_0(a)-a))+\tau(T_0(a^2)-a^2)
 \le3\eta.
 \label{u:eq:averaged-three}
\end{align}
In the last inequality each trace term is bounded by the norm of its
argument, and $\norm{a}\le1$.

Apply Lemma~\ref{u:lem:homomorphism-replacement}, with the same $\eta$,
to each $\phi_i$, obtaining $\rho_i:F_i\to D$.  Since $\psi_i(a)$ is a
contraction, Equations~\eqref{u:eq:homomorphism-replacement} and
\eqref{u:eq:averaged-unit-defect} give
\[
 \sum_i\beta_i
 \norm{T_i(a)-\rho_i\psi_i(a)}_{2,\tau}^{2}
 \le4\sum_i\beta_i(1-\tau(h_i))+2\eta\le6\eta.
\]
A second squared triangle inequality and
Equation~\eqref{u:eq:averaged-three} yield
\begin{equation}
 \sum_i\beta_i\norm{a-\rho_i\psi_i(a)}_{2,\tau}^{2}
 \le2(3\eta)+2(6\eta)=18\eta.
 \label{u:eq:averaged-eighteen}
\end{equation}

Choose the $r_i$ successively.  At step $i$ apply
Proposition~\ref{u:prop:centered-symmetry} with
\[
 E_i=\rho_i(F_i),\qquad S_i=C^*(A,r_1,\ldots,r_{i-1}).
\]
The algebra $S_i$ is unital and norm separable, while $E_i$ is finite
dimensional and may be nonunital.  We obtain
\begin{equation}
 r_i=r_i^*,\quad r_i^2=1,\quad r_i\in E_i',\quad
 \tau(sr_i)=0\quad(s\in S_i,\ \tau\in\Lambda).
 \label{u:eq:successive-centering}
\end{equation}
Because $r_i$ commutes with $\rho_i\psi_i(a)$ and multiplication by a
unitary preserves the tracial $2$-seminorm,
\[
 \norm{[r_i,a]}_{2,\tau}
 \le2\norm{a-\rho_i\psi_i(a)}_{2,\tau}.
\]
Thus Equation~\eqref{u:eq:averaged-eighteen} gives
\begin{equation}
 \sum_i\beta_i\norm{[r_i,a]}_{2,\tau}^2\le72\eta
 \quad(a\in E).
 \label{u:eq:averaged-seventytwo}
\end{equation}
Every weighted estimate has been obtained at the same fixed $\tau$;
their constants are independent of that trace.

Set $z=\sum_i\sqrt{\beta_i}\,r_i$.  We spell out the cancellation that
converts Equation~\eqref{u:eq:averaged-seventytwo} into a bound for $z$.
If a word contains $r_j$ exactly once and all its other letters belong
to $S_j$, cyclicity writes its trace as $\tau(sr_j)$ for some $s\in S_j$.
It vanishes by Equation~\eqref{u:eq:successive-centering}.  For $i<j$,
each of the four terms in
$[r_i,a]^*[r_j,a]$ has this form, since $a,r_i\in S_j$.  The reverse
mixed term is its adjoint and also has zero trace.  Therefore
\[
 \norm{[z,a]}_{2,\tau}^2
 =\sum_i\beta_i\norm{[r_i,a]}_{2,\tau}^2\le72\eta.
\]
For any $b\in A$, the same argument gives $\tau(br_i)=0$ and
$\tau(br_ir_j)=0$ for $i\ne j$.  The diagonal terms use $r_i^2=1$ and
$\sum_i\beta_i=1$, proving
$\tau(bz)=0$ and $\tau(bz^2)=\tau(b)$.  These are coefficient identities,
and do not require the different $r_i$ to commute.

For the fourth moment write
$z_i=\sum_{j\le i}\sqrt{\beta_j}\,r_j$, with $z_0=0$.  At step $i$ put
$X=z_{i-1}$, $r=r_i$, and $c=\sqrt{\beta_i}$.  The full noncommutative
expansion, grouped only after applying the trace, is
\begin{align*}
 \tau((X+cr)^4)
 ={}&\tau(X^4)+4c\tau(X^3r)\\
 &+c^2\bigl(4\tau(X^2)+2\tau(XrXr)\bigr)
   +4c^3\tau(Xr)+c^4.
\end{align*}
Indeed, the four words with one $r$ are cyclic permutations of $X^3r$.
The four words with three $r$ reduce, by cyclicity and $r^2=1$, to
$Xr$.  Of the six words with two $r$'s, four have trace $\tau(X^2)$,
and two are alternating.  The terms with one or three $r$ vanish by
Equation~\eqref{u:eq:successive-centering}.  Tracial Cauchy--Schwarz gives
\[
 |\tau(XrXr)|
 \le\norm{Xr}_{2,\tau}\norm{rX}_{2,\tau}
 =\tau(X^2).
\]
The already proved second-moment cancellation, applied to the partial
sum, gives $\tau(X^2)=\sum_{j<i}\beta_j$.  Consequently
\[
 \tau(z_i^4)\le\tau(z_{i-1}^4)
       +6\beta_i\sum_{j<i}\beta_j+\beta_i^2.
\]
Summing over $i$ proves the claimed bound, including its sharper form:
\begin{equation}
 \tau(z^4)\le6\sum_{j<i}\beta_i\beta_j+\sum_i\beta_i^2
   =3\Bigl(\sum_i\beta_i\Bigr)^2-2\sum_i\beta_i^2
   =3-2\sum_i\beta_i^2\le3.
 \label{u:eq:averaged-fourth}
\end{equation}
Finally, $\norm{z}\le\sum_i\sqrt{\beta_i}\le\sqrt M$, which explains
why this construction alone does not give a common norm bound as the
approximation data vary.
\end{proof}

\subsection{Clipping and the local assertion in \texorpdfstring{$A$}{A}}

\begin{lemma}[Tracial clipping]
\label{u:lem:clipping}
Let $z=z^*\in D$, let $R>0$, and set
\[
 g_R(s)=\max(-R,\min(s,R)),\qquad d=g_R(z).
\]
For every $a\in D$ and $\tau\in\Lambda$,
\begin{equation}
 \norm{[d,a]}_{2,\tau}\le\norm{[z,a]}_{2,\tau}.
 \label{u:eq:clipping-commutator}
\end{equation}
Moreover, for every contraction $b\in D$,
\begin{equation}
 \begin{aligned}
  |\tau(b(d-z))|&\le R^{-3}\tau(z^4),\\
  |\tau(b(d^2-z^2))|&\le R^{-2}\tau(z^4),\\
  \tau(d^4)&\le\tau(z^4).
 \end{aligned}
 \label{u:eq:clipping-tails}
\end{equation}
In particular, if $z$ is obtained from
Proposition~\ref{u:prop:averaged-observable}, then $\norm{d}\le R$ and,
for every $\tau\in\Lambda$,
\begin{equation}
 \begin{aligned}
  \norm{[d,a]}_{2,\tau}&\le\sqrt{72\eta} &&(a\in E),\\
  |\tau(bd)|&\le3/R^3 &&(b\in A,\ \norm{b}\le1),\\
  |\tau(b(d^2-1))|&\le3/R^2 &&(b\in A,\ \norm{b}\le1),\\
  \tau(d^4)&\le3
 \end{aligned}
 \label{u:eq:clipped-moments}
\end{equation}
\end{lemma}

\begin{proof}
Fix $\tau$ and form the Hilbert space $L^2(D,\tau)$ by quotienting out
the zero vectors of the tracial seminorm and completing.  Left and
right multiplication by the bounded selfadjoint $z$ define commuting
bounded selfadjoint operators $L_z$ and $R_z$ on this space.  Continuous
functional calculus gives $g_R(L_z)=L_d$ and $g_R(R_z)=R_d$.
Let $E_z$ be their joint spectral measure, and let $\mu_a$ be the finite
positive measure $\mu_a(U)=\langle E_z(U)a,a\rangle$.  The scalar
$1$-Lipschitz inequality for $g_R$ gives
\begin{align*}
 \norm{[d,a]}_{2,\tau}^2
 &=\int_{\R^2}|g_R(s)-g_R(t)|^2\,d\mu_a(s,t)\\
 &\le\int_{\R^2}|s-t|^2\,d\mu_a(s,t)
  =\norm{[z,a]}_{2,\tau}^2.
\end{align*}
This proves Equation~\eqref{u:eq:clipping-commutator} in tracial $L^2$.

For real $s$, the scalar inequalities
\[
 |s-g_R(s)|\le\frac{|s|^4}{R^3},\qquad
 |s^2-g_R(s)^2|\le\frac{|s|^4}{R^2},\qquad
 |g_R(s)|^4\le|s|^4
\]
follow by considering $|s|\le R$ and $|s|>R$.  Functional calculus
therefore bounds the traces of $|z-d|$ and $|z^2-d^2|$ by the first two
right-hand sides of Equation~\eqref{u:eq:clipping-tails}.
For completeness, the trace inequality that permits arbitrary
contractive coefficients is
\begin{equation}
 |\tau(bx)|\le\tau(|x|)\quad(\norm{b}\le1).
 \label{u:eq:trace-coefficient-bound}
\end{equation}
To see it, take the polar decomposition $x=u|x|$ in $D^{**}$ and use the
normal tracial extension of $\tau$.  With $y=|x|^{1/2}$ and $c=bu$,
traciality and Cauchy--Schwarz give
\[
 |\tau(bx)|=|\tau(ycy)|
 \le\tau(y^2)^{1/2}\tau(yc^*cy)^{1/2}
 \le\tau(y^2)=\tau(|x|).
\]
This proves Equation~\eqref{u:eq:clipping-tails} without a commutation
assumption on $b$.  Combining it with
Equation~\eqref{u:eq:averaged-observable} proves
Equation~\eqref{u:eq:clipped-moments}.
\end{proof}

\begin{proposition}[Local bounded moments in $A$]
\label{u:prop:local-moments}
Fix $R\ge1$.  For every finite set $E\subset A$ of selfadjoint
contractions, every finite set $H\subset A$ of contractions, and every
$\sigma>0$, there is a selfadjoint $d\in A$ with $\norm{d}\le R$ such
that, simultaneously,
\begin{equation}
 \begin{aligned}
  \norm{[d,a]}_{2,T}&<\sigma &&(a\in E),\\
  \sup_{\tau\in T}|\tau(bd)|&\le3/R^3+\sigma &&(b\in H),\\
  \sup_{\tau\in T}|\tau(b(d^2-1))|&\le3/R^2+\sigma &&(b\in H),\\
  \sup_{\tau\in T}\tau(d^4)&\le3+\sigma.
 \end{aligned}
 \label{u:eq:local-moments}
\end{equation}
\end{proposition}

\begin{proof}
Choose $0<\eta<1$ so that $\sqrt{72\eta}<\sigma/2$.
Proposition~\ref{u:prop:averaged-observable} and Lemma~\ref{u:lem:clipping}
give $\bar d=\bar d^*\in D$ with $\norm{\bar d}\le R$ satisfying
Equation~\eqref{u:eq:clipped-moments}.  Choose a selfadjoint lift of
$\bar d$ to $\ell^\infty(\N,A)$ and apply $g_R$ at every coordinate.
The resulting lift $(d_l)$ remains a lift of $\bar d$ and satisfies
$d_l=d_l^*$ and $\norm{d_l}\le R$ for every $l$.

By Equation~\eqref{u:eq:limit-trace-norm} and the trace-test identity in
Lemma~\ref{u:lem:tracial-diagonal},
\begin{align*}
 \lim_{l\to\omega}\norm{[d_l,a]}_{2,T}
 &=\norm{[\bar d,a]}_{2,\Lambda}<\sigma/2,\\
 \lim_{l\to\omega}\sup_{\tau\in T}|\tau(bd_l)|
 &=\sup_{\tau\in\Lambda}|\tau(b\bar d)|\le3/R^3,\\
 \lim_{l\to\omega}\sup_{\tau\in T}|\tau(b(d_l^2-1))|
 &=\sup_{\tau\in\Lambda}|\tau(b(\bar d^2-1))|\le3/R^2,\\
 \lim_{l\to\omega}\sup_{\tau\in T}\tau(d_l^4)
 &=\sup_{\tau\in\Lambda}\tau(\bar d^4)\le3.
\end{align*}
There are only finitely many tests, indexed by $E$, $H$, and the one
fourth-moment test.  For each test the desired bound in
Equation~\eqref{u:eq:local-moments} holds on a set in $\omega$, using the
strict margin $\sigma/2$ for commutators and the positive margin
$\sigma$ for moments.  Their finite intersection belongs to $\omega$
and is nonempty.  Any coordinate in it supplies the required $d$.
\end{proof}

For fixed $R$ the two coefficient errors retained in
Equation~\eqref{u:eq:local-moments} are $3/R^3$ and $3/R^2$; only the
additional error $\sigma$ is arbitrarily small.  The norm bound $R$ is
independent of the nuclear approximation and its number of summands.
These are the precise features needed for the passage to separable
coefficient algebras in the next section.
